\subsection{Example 2}\label{sec:interval-example2}

\subsubsection{Problem description}\label{subsec:interval-example2-problem}

A finance team is estimating the average monthly return (\%) of a diversified investment portfolio based on a very large sample of return observations. Because the dataset is large, the t-student distribution can be approximated by the normal distribution, even though the population variance is unknown. The team wants to build a 95\% confidence interval for the mean return to guide risk and allocation decisions.

Given sample summary:

\begin{itemize}
	\item Sample size (\(n = 400\)).
	\item Sample mean (\(\bar{x} = 1.2\%\)).
	\item Sample standard deviation (\(s = 2.5\%\)).
\end{itemize}

Estimate the 95\% confidence interval for the population mean monthly return (\(\mu\)), and explain why the normal approximation is justified.


\subsubsection{(C9) Estimate the mean with a large-sample t approximation}\label{subsubsec:interval-example2-c9}

Because the population variance is unknown, the exact approach would use the t distribution with \(n - 1\) degrees of freedom. However, since \(n = 400\) is very large, the t distribution is almost identical to the normal distribution. Therefore we approximate using \(z_{0.025} = 1.96\).

Compute the margin of error:

\[
E \approx z_{\alpha/2} \cdot \frac{s}{\sqrt{n}}
= 1.96 \cdot \frac{2.5}{\sqrt{400}}
= 1.96 \cdot \frac{2.5}{20}
= 1.96 \cdot 0.125
= 0.245
\]

So the confidence interval is:

\[
\mu \in 1.2 \pm 0.245 = (0.955, 1.445)
\]

This interval estimates the plausible range for the true mean monthly return.

\subsubsection{(C4) Interpret the confidence interval for a finance decision}\label{subsubsec:interval-example2-c4}

The interval suggests the portfolio's true mean monthly return is likely between about 0.96\% and 1.45\%. For allocation decisions, the finance team can use the lower bound as a conservative expected return when stress-testing plans and use the midpoint as a baseline forecast. The range also helps compare this portfolio to alternatives with different risk profiles.

\subsubsection{(C10) Conclude about a statistical parameter in finance and economy}\label{subsubsec:interval-example2-c10}

Based on the large sample, the team can conclude that the population mean monthly return is about 1.2\%, with a 95\% confidence interval from roughly 0.96\% to 1.45\%. Even though the variance is unknown, the large-sample approximation supports using the normal critical value to obtain a reliable interval for planning.

\vspace{6pt}
\noindent\hyperref[table-of-contents]{Back to Top}
